\subsection{\texorpdfstring{Functions of class \(C^{r}\) ( \(r \neq \omega\) )}{Functions of class C\^{}\{r\} (r ≠ ω)}}

2.3.1. Let U be an open subset of E and f a mapping from U into F. One defines the relation “f is of class \(C^{r}\) ” (for \(r \in N\) ) by induction on r in the following manner:

\begin{enumerate}
\def\labelenumi{\arabic{enumi})}
\item
  \(f\) is of class \(\mathbf{C}^{0}\) if and only if it is continuous;
\item
  if \(r\) is an integer \(\geqslant 1\) , the function \(f\) is of class \(\mathbf{C}^{r}\) if and only if it is differentiable at every point of \(\mathbf{U}\) and if the derivative \(Df\) from \(\mathbf{U}\) to \(\mathcal{L}(\mathbf{E};\mathbf{F})\) is of class \(\mathbf{C}^{r-1}\) .
\end{enumerate}

The functions of class C \(^{r}\) are also called r times continuously differentiable functions.

One says that \(f\) is of class \(\mathbf{C}^{\infty}\) (or infinitely differentiable) if it is of class \(\mathbf{C}^{r}\) for every integer \(r\) .

If \(f\) is of class \(\mathbf{C}^{r}\) in \(\mathbf{U}\) , then \(f\) is \(p\) times differentiable for every integer \(p \leqslant r\) and the function \(\mathbf{D}^{p}f\) is of class \(\mathbf{C}^{r-p}\) .

2.3.2. The mappings of class \(C^{r}\) from an open subset U of E into F form a vector subspace \(\mathcal{C}^{r}(U;F)\) of the space of all mappings from U into F. One has \(\mathcal{C}^{s}(U;F)\subset\mathcal{C}^{r}(U;F)\) for \(s\geqslant r\) .

2.3.3. For a function \(f\) to be of class \(\mathbf{C}^{1}\) in an open subset U of E, it is necessary and sufficient that it be strictly differentiable at every point of U.

If E is a product of normed spaces \(E_{i}\), a map \(f\) from an open set V of E into F is of class \(C^{r}\) if and only if \(f\) possesses iterated partial derivatives \(D_{i_{1}}\ldots D_{i_{m}}f\) continuous for every integer \(m \leqslant r\) .

2.3.4. Let G be a normed space and let U be an open subset of E. Let V be an open subset of G, \(g \in \mathcal{C}^{r}(U; G)\) and \(f \in \mathcal{C}^{r}(V; F)\) . If \(g(U) \subset V\) , the map \(f \circ g\) of U into F is of class \(C^{r}\) .

Let \(\mathbf{F}_1\) and \(\mathbf{F}_2\) be two separated locally convex spaces and let \(u\) be a bilinear mapping from \(\mathbf{F}_1 \times \mathbf{F}_2\) to \(\mathbf{F}\) , hypocontinuous with respect to the set of bounded subsets of \(\mathbf{F}_1\) (resp. \(\mathbf{F}_2\) ) (Topological Vector Spaces, ch. III, § 4, no. 2). Let \(\mathbf{U}\) be an open subset of \(\mathbf{E}\) and \(f_i \in \mathcal{C}^r(\mathbf{U};\mathbf{F}_i)\) (for \(i = 1,2\) ). Then the function \(u(f_1,f_2)\) belongs to \(\mathcal{C}^r(\mathbf{U};\mathbf{F})\) . If \(\mathbf{E}\) is finite-dimensional, it suffices to assume that \(u\) satisfies condition (SC) of no. 2.1.3.

2.3.5. If E is a product of normed spaces \(E_{i}\) ( \(1 \leqslant i \leqslant n\) ) and if f is a continuous n-linear mapping from E into F, then f is of class \(C^{\infty}\) and one has \(D^{p}f = 0\) for \(p \geqslant n + 1\) .

2.3.6. Suppose that E and F are Banach spaces. Let f be a function of class \(C^{r}\) (with \(r \geqslant 1\) ) defined in a neighborhood of a point \(x_{0}\) of E and with values in F. Let \(y_{0} = f(x_{0})\) and suppose that \(Df(x_{0})\) is an isomorphism from E onto F. Then f induces a homeomorphism g from a neighborhood of \(x_{0}\) onto a neighborhood of \(y_{0}\) (1.5), and the inverse mapping of g is of class \(C^{r}\) in a neighborhood of \(y_{0}\) .

